\chapter{统一实验设置、指标体系与复现细节}

\section{本章写作目的}
前几章分别介绍了细胞核实例分割、零样本核检测、图像提示目标检测以及肺鳞癌数字病理分析四个研究工作。由于这些工作发表或投稿于不同会议和期刊，其原始论文在实验组织方式、指标命名、图表排版和任务侧重点上并不完全一致。为了使整篇学位论文形成统一的技术叙事，本章对数据集、实验任务、评价指标、训练策略和复现细节进行集中说明。

这一章的作用有两点。第一，它为读者提供统一背景，避免在不同章节之间反复解释相同概念。第二，它把作者在多个工作中逐步形成的实验规范系统化，包括固定支持样本、跨域数据划分、直接测试与目标域微调的区别、以及检测和分割指标的对应关系。这些内容虽然在单篇小论文中可能只是实现细节，但在博士或硕士学位论文中十分重要，因为它们体现了研究工作的连续性和可复现性。

\section{数据集与任务口径}
\subsection{细胞核检测与分割数据集}
本文涉及的细胞核任务主要基于多种公开或合作构建的病理图像数据集，包括 MoNuSAC、CoNSeP、Lizard、PanNuke 以及肾透明细胞癌相关图像。不同数据集在组织类型、染色风格、类别定义和实例密度上存在明显差异。例如，MoNuSAC 更强调多器官场景下的核类别多样性，CoNSeP 更强调结直肠癌组织中的细胞核亚型，Lizard 则具有更大规模和更复杂的组织背景。正是这些差异，使它们非常适合用于检验跨域检测和分割方法的泛化能力。

在学位论文中，需要明确区分“源域训练集”和“目标域测试集”。源域训练集用于学习通用表示或模型初始能力，目标域则用于检验模型面对新数据分布时的迁移表现。对于 few-shot 或 image-prompt 设定，还需要额外定义目标域支持集，即从目标域训练图像中抽取少量图像作为提示或适配信息。本文在后续实验中尽量固定不同方法使用的支持图像，从而保证比较结果不受随机抽样差异影响。

\subsection{肺癌数字病理数据}
第六章使用的肺鳞癌数字病理数据与前几章的 patch 级核检测分割数据不同。该工作更接近临床队列研究，数据单位通常是患者、切片或全切片图像。其重点并不是单张图像中的检测精度，而是 AI 量化指标与病理缓解程度、免疫细胞密度、无病生存和总生存之间的关系。因此，该章实验需要同时考虑图像处理流程和临床统计流程。

这类数据的特殊性在于，模型输出并不是最终结论，而是临床指标构建的中间层。例如，残余可存活肿瘤比例需要由组织区域识别结果推导，肿瘤周围免疫细胞密度需要由细胞检测结果和空间边界关系共同决定。相比检测或分割 benchmark，这类研究更加重视量化指标的可解释性、稳定性和临床相关性。

\section{跨域任务设置}
\subsection{源域到目标域的迁移}
本文多个工作都围绕跨域医学图像任务展开。典型任务包括 MoNuSAC 到 CoNSeP、CoNSeP 到 MoNuSAC、CoNSeP 到 Lizard、MoNuSAC 到 Lizard、PanNuke 到 MoNuSAC、CoNSeP 到 PanNuke、MoNuSAC 到 CCRCC 以及 CoNSeP 到 CCRCC。每个任务都可以记为
\[
D_{\text{source}} \rightarrow D_{\text{target}},
\]
其中源域提供模型训练或预训练，目标域用于测试模型的迁移能力。

跨域任务比同域任务更接近真实病理应用。现实中，模型往往不可能在每个医院、每种组织和每种染色流程上都重新获取大量标注。一个有实际价值的医学视觉模型，必须能够在源域学到可迁移的结构知识，并通过少量提示或少量支持样本快速适配目标域。本文第三章和第五章提出的提示学习方法，正是围绕这一需求展开。

\subsection{少样本支持集设置}
对于少样本实验，本文通常使用 1-shot、5-shot、10-shot 和 50-shot 四种设置。这里的 shot 指每个类别从目标域训练集中选取的支持图像数量。为了计算平均值和标准差，实验通常使用 3 个随机种子生成三组不同支持集。需要强调的是，支持图像的选择必须在不同方法之间保持一致，否则方法差异可能被数据抽样差异掩盖。

在图像提示方法中，支持集有两种不同用途。第一，它可以作为目标域微调数据，用于更新模型部分参数。第二，它也可以只作为提示输入，在不更新模型权重的情况下为测试图像提供目标域上下文。前者通常可能获得更高性能，但需要额外训练成本；后者更加轻量，适合快速部署和公平比较。本文后续主表和消融实验中，需要明确区分这两类设置。

\section{直接测试与目标域微调}
\subsection{直接测试模式}
直接测试模式指模型加载源域训练好的权重后，不在目标域执行梯度更新，而是直接利用目标域支持图像构造提示并进行推理。该模式可以写作
\[
f_{\theta_{\text{source}}}(x_{\text{test}} \mid \mathcal{S}_{\text{target}}),
\]
其中 $\theta_{\text{source}}$ 表示源域训练得到的模型参数，$\mathcal{S}_{\text{target}}$ 表示目标域少量支持图像。这个设定能够更纯粹地检验提示机制本身的作用，因为目标域性能来自源域知识和支持提示，而不是来自目标域梯度优化。

直接测试模式的优点是成本低、速度快、复现简单，也更适合说明“模型是否真的具备提示驱动的泛化能力”。其缺点是性能通常低于目标域微调，因为模型参数没有机会根据目标域图像进行适配。在学位论文中，可以把这一模式作为低成本泛化能力的证据，而不是简单与充分微调结果混为一谈。

\subsection{目标域微调模式}
目标域微调模式允许模型使用目标域少量支持图像更新部分参数。根据冻结策略不同，可以只微调检测头和分割头，也可以微调更大范围的模型模块。该模式通常更接近“追求目标域最优性能”的实验设置，但其训练时间和计算资源消耗更高，也更容易受到支持集数量和随机种子的影响。

本文在组织实验时需要同时保留这两类结果。直接测试结果用于说明方法的即插即用能力，目标域微调结果用于说明方法在有限标注条件下的上限潜力。二者回答的问题不同，因此在主表或消融表中应当清楚标注，避免读者误以为所有数值来自同一种实验条件。

\section{评价指标体系}
\subsection{检测指标}
检测任务主要使用平均精度和平均召回率进行评价。平均精度通常记为 AP 或 mAP，用于衡量模型预测框在不同置信度阈值和重叠阈值下的准确性。平均召回率通常记为 AR，用于衡量模型是否能够尽可能找全真实目标。在核检测任务中，召回率尤其重要，因为漏检会直接影响后续计数、空间分布分析和微环境量化。

不过，单独依赖 AP 或 AR 也存在局限。AP 更偏向排序质量和定位精度，AR 更偏向目标覆盖能力。对于密集核目标，模型可能出现较多重叠框或碎片框，使 AP 与实际可用性之间存在差距。因此，本文在讨论检测结果时，通常结合定性可视化和后续分割指标共同判断模型行为。

\subsection{实例分割指标}
实例分割部分使用 Dice、Hausdorff Distance、AJI、对象级 Dice、PQ 和 F1 等指标。Dice 衡量整体前景重叠程度，适合评估模型是否大体找到了目标区域。Hausdorff Distance 反映边界误差，数值越低通常表示边界越接近真实标注。AJI 强调实例级聚合交并比，适合评价核实例分割中的整体实例质量。对象级 Dice 在匹配实例后计算平均重叠程度，更关注单个细胞核对象是否被正确分割。

PQ 和 F1 更强调实例检测与分割的联合质量。特别是在密集核场景中，模型可能获得较高前景 Dice，但由于多个细胞核粘连为一个大实例，PQ 和 F1 仍然较低。这类现象在本文实验中非常常见，因此结果分析不能只看单一 Dice 值，而应同时观察对象级指标。

\subsection{临床统计指标}
第六章涉及的临床统计指标包括组间差异、相关性分析、ICC 一致性、风险比以及生存曲线等。这些指标与检测和分割指标不同，其目标不是评价单张图像预测结果，而是判断 AI 量化指标是否与病理评估和患者结局存在稳定关系。对于这类实验，模型性能的最终价值体现在能否形成可解释的临床分层，而不是单纯在图像层面取得高分。

\section{图表整理原则}
由于本文由多项已发表或在投工作整合而成，原始图表大多以英文形式出现在小论文中。在大论文中使用这些图表时，需要进行三类整理。第一，图题和表题应改为中文，并明确指出图表服务于哪个实验问题。第二，图中英文标注可以保留原始缩写，但正文中应给出中文解释。第三，若图表过于密集，可将主体结论放在正文，完整原图放入附录。

这一整理过程并不是简单翻译，而是重新建立中文学位论文的叙事逻辑。例如，会议论文中的一张主结果表可能只是为了证明 SOTA，而在学位论文中，它还应承担说明数据集设置、任务难度和方法贡献边界的作用。因此，图表进入大论文后，需要配套增加更充分的文字分析。

\section{复现与工程规范}
本文的实验复现涉及多个层面，包括数据集路径、随机种子、支持集抽样、模型权重、配置文件和评价脚本。为了保证结果可追溯，每个实验应尽量记录运行命令、输出目录、预测 JSON、模型权重和日志文件。对于 few-shot 实验，还应保存支持集 manifest，确保后续不同方法能够使用同一批支持图像。

在工程管理上，直接测试、目标域微调和源域训练应使用不同输出目录，避免结果互相覆盖。对于长时间训练任务，还需要区分中间 checkpoint、最终模型和推理结果。由于病理图像实验通常会产生大量预测缓存，论文整理阶段可以保留必要的 mask JSON、配置文件和日志，而删除体积过大的中间缓存，以降低工程目录体积。

\section{本章小结}
本章对全文实验设置进行了统一整理，重点说明了数据集口径、跨域任务、少样本支持集、直接测试与目标域微调、评价指标以及复现规范。通过这一章，读者可以更清楚地理解前几章实验之间的共同基础，也能更容易判断不同表格中结果数值的实际含义。对于后续正式送审版本，本章还可以继续扩展为完整的“实验设置与实现细节”章节，从而进一步增强论文的完整性和可复现性。
